\Qhead{4.10}{Взаимодействие света со звуком. Акустооптические устройства}{Взаимодействия света со звуком. Дифракция Брэгга и Рамана-Ната. Принципы работы устройств акустооптики (модуляторы и дефлекторы света, преобразователи свет-сигнал, акустооптические фильтры), анализаторы спектра и корреляторы.}
\Subfile{4.10.1}{sections/04_generation/q10/s01}
\Subfile{4.10.2}{sections/04_generation/q10/s02}

\recap{Звук создаёт движущуюся фотоупругую решётку. Толщина и синхронизм выбирают много порядков Рамана–Ната или двухволновый Брэгг. Зависимость эффективности и угла от звука позволяет модулировать, сканировать, фильтровать и обрабатывать сигналы.}
\board{Векторный треугольник, $Q_{\rm KN}$, две связанные оптические амплитуды; $N\sim\Delta f\tau$; линза спектроанализатора.}
\oral{Почему формула Брэгга с одним показателем преломления может не работать при смене оптической поляризации в анизотропном кристалле?}
\bait{Можно оставить компромисс апертуры и быстродействия: та же акустическая задержка повышает число разрешимых частот и замедляет переключение.}
\sources{Преобразование изображения в сигнал: Балакший и соавт.~\cite{balakshylightsignal}, с.~743, описание работы АРУС (первичный текст); Университетские слайды «Основы оптоэлектроники»~\cite{optoelectronics}, раздел управления светом; Warde~\cite{wardeao}, лабораторные материалы по АО-модуляции. Фотоупругая связь и векторные условия используются в явных соглашениях; полное пособие Балакшия, Парыгина, Чиркова из программы не получено.}
